\documentclass[12pt]{article}
\usepackage[utf8]{inputenc}
\usepackage[T1]{fontenc}
\usepackage{hyperref}
\title{Patla-Chicontla Totonac Language Resources for K-12 Curriculum Development}
\author{Compiled by Claude Sonnet 5 (claude-sonnet-5) for Andr\'as Kornai}
\date{}
\begin{document}
\maketitle

\section{Introduction}
The identification block for this language was incomplete: the family field read ``Bookkeeping'', which is not a language family. Sources found describe Patla-Chicontla Totonac as a Totonacan language of the Sierra Norte of Puebla, spoken in Patla, Chicontla and nearby villages. Some sources treat it as the same as Upper Necaxa Totonac, which has the code tku. I could not confirm the relationship between the codes tot and tku, and it should be checked against an authoritative registry. The material below may therefore belong to a neighbouring variety. Reported speaker numbers are about 3,400, mostly adults.

\section{Primary Educational Resources}
None identified.

\section{Linguistic and Reference Materials}
Search results included pages on a University of Alberta Totonac site. I did not see their content, so I cannot say whether they are dictionaries or grammars. They are listed only as leads.
\begin{itemize}
\item \url{https://www.artsrn.ualberta.ca/totonaco2/?p=105}
\item \url{https://www.artsrn.ualberta.ca/totonaco2/?p=146}
\item \url{https://www.artsrn.ualberta.ca/totonaco2/?p=123}
\item Omniglot, Upper Necaxa Totonac: \url{https://omniglot.com/writing/uppernecaxatotonac.htm}
\end{itemize}

\section{Religious and Community Resources}
\begin{itemize}
\item Joshua Project profile: \url{https://joshuaproject.net/index.php/people_groups/14339}
\item Scripture Earth listing, under code tku: \url{https://scriptureearth.org/00eng.php?iso=tku}
\end{itemize}

\section{Digital and Audio Resources}
No specific audio resources were confirmed.

\section{Curriculum Development Context}
No school programme was found. Little confirmed material exists, and the code and identity ambiguity should be resolved first.
\end{document}